\section{System Model and Problem Formulation}
\label{sec:system_model}

\subsection{Distributed Learning Setup}

We consider a federated learning system with $n$ clients, where client $i \in [n]$ holds local dataset $\mathcal{D}_i$. The global objective is:
\begin{equation}
F(w) = \frac{1}{n}\sum_{i=1}^n F_i(w),\quad F_i(w) = \mathbb{E}_{z \sim \mathcal{D}_i}[\ell(w;z)]
\end{equation}

At each round $t \in [T]$, the aggregator broadcasts $w^t$. Each honest client $i$ computes:
$g_i^t = \nabla F_i(w^t) + \xi_i^t$,
where $\xi_i^t$ is stochastic noise with coordinate-wise variance $\mathbb{E}[(\xi_i^t)_j^2] \leq \sigma^2$.

\begin{figure}[t]
\centering
\includegraphics[width=\textwidth]{flowdiag.jpeg}
\caption{System architecture. Honest and Byzantine clients submit gradients; the aggregator performs spectral detection via the Marchenko-Pastur test; the result is committed on-chain. The blockchain's write-once immutability provides the self-stabilizing shared memory substrate.}
\label{fig:architecture}
\end{figure}

\subsection{Formal Distributed System Model}
\label{sec:dist_model}

We model the system as a synchronous message-passing distributed system~\cite{byzantine_generals}.

\smallskip\noindent\textbf{Processes.} $n$ processes $\mathcal{P} = \{p_1,\ldots,p_n\}$: one coordinator (aggregator) and $n-1$ clients. Each process has local state $w^t \in \mathbb{R}^d$.

\smallskip\noindent\textbf{Execution.} Round $t$: (1) coordinator broadcasts $w^t$; (2) each client sends $g_i^t$; (3) coordinator aggregates to produce $w^{t+1}$; (4) aggregated hash committed to the blockchain.

\smallskip\noindent\textbf{Configurations.} $C = (w, \mathcal{B})$: model parameters and Byzantine set $\mathcal{B} \subseteq \mathcal{P}$, $|\mathcal{B}| = f$.

\subsection{Threat Model and Assumptions}
\label{sec:threat_model}

\noindent\textbf{Assumption~1 (Byzantine Fault Model).}
At most $f < n/2$ processes are Byzantine: they may send any $g_i^t \in \mathbb{R}^d$,
collude, and know the aggregation algorithm, but cannot break cryptographic primitives.

\noindent\textbf{Assumption~2 (Honest Gradient Variance).}
Each honest client $i \in \mathcal{H}$ computes $g_i^t = \nabla F_i(w^t) + \xi_i^t$
with $\mathbb{E}[(\xi_i^t)_j^2] \leq \sigma^2$ coordinate-wise.

\noindent\textbf{Assumption~3 (Phase Transition Condition).}
Let $\beta = f/n \in (0, \tfrac{1}{2})$ denote the Byzantine fraction.
We require $\sigma^2\beta^2 < 0.25$, equivalently $\sigma\beta < \tfrac{1}{2}$.
This necessary condition ensures Byzantine gradients create
detectable spectral anomalies (Theorem~\ref{thm:phase_impossibility});
the matching impossibility shows the condition cannot be relaxed.

\noindent\textbf{Assumption~4 (Blockchain Liveness).}
The blockchain finalizes honest transactions with at most $\lfloor(m-1)/3\rfloor$ faulty
validators out of $m$ (PBFT assumption~\cite{bft_consensus}).

\noindent\textbf{Assumption~5 (Transient Faults).}
Byzantine behavior may occur at any round including round~0.
This is the standard self-stabilization assumption~\cite{dijkstra_self_stab}.

\smallskip\noindent\textbf{Attack taxonomy.}
Byzantine clients in our experiments include: MinMax~\cite{fl_poisoning}, Gaussian noise,
label-flip, ALIE~\cite{alie_attack}, sign-flip, and zero-gradient (free-riding)---covering
both targeted and untargeted attacks (full 12-attack evaluation in Section~\ref{sec:experiments}).

\subsection{Protocol Properties}
\label{sec:properties}

\noindent\textbf{Safety.} $\mathbb{E}[\|\hat{g}^t - \bar{g}^t\|^2] \leq O(\sigma^2\beta^2)$ in every round, where $\beta = f/n$.

\noindent\textbf{Liveness.} $\min_{t \in [T]}\mathbb{E}[\|\nabla F(w^t)\|^2] \leq O(\sigma\beta/\sqrt{T} + \beta^2/T)$
provided $f < n/2$ and $\sigma^2\beta^2 < 0.25$.

\noindent\textbf{Practical Stabilization.} Starting from any $w^0 \in \mathbb{R}^d$, the protocol reaches
a legitimate configuration (Definition~\ref{def:legit}) within
$T^* = O(\sigma^2\beta^2/\varepsilon^2 + \beta^2/\varepsilon)$ rounds (Theorem~\ref{thm:self_stab}).
